% !TEX root = ../main.tex

\begin{acknowledgements}

我首先要感谢我的导师李嘉禹教授，我很荣幸也很幸运能成为李老师的学生。我很感激李老师在我初学复几何时给予我的耐心和鼓励。在博士中期，李老师向我推荐了抛物丛这个有趣且有深度的课题。这篇文章很多的想法都是建立在李老师早期的研究工作之上。我也很感激李老师在博士期间给予我的学术自由，他从来没有限制我的兴趣，包容我的任性，给我足够大的空间去探索复几何的世界，这对我尤为重要。在学术上，李老师的指导言简意赅，他总是能抓到最关键的信息，这给了我很大的帮助。在为人处事上，李老师谦逊温和，是一个真正的绅士，为我树立了榜样。

当然我要感谢我的父母。我感谢我的父亲蒋世雨先生对我选择基础数学研究的理解和支持，这是我的精神支柱。我也感谢我的母亲蔡伟玲女士在生活上给予我无微不至的关怀。感谢他们在我未如愿时给予我的信心，在风雨里带给我的温暖。

我还要感谢Takuro Mochizuki教授对我的帮助。在博士期间，我与他有过多次的邮件交流，他每次都能精准详细地解答我的疑惑。我感谢他能够把我当成一个同行平等地交流并作为长辈给予我合理的建议。同时我感谢他为数学所做出的贡献。他对数学几十年如一日的真诚和热情是我要学习的。

我还要感谢张世宇学长组织的Nonabelian Hodge对应主题的讨论班以及潘长鹏师兄的参与，和他们的交流让我受益良多。

最后我要感谢我的师弟张一剑。感谢他带着热情和认真的学习态度和我共同成长。
\end{acknowledgements}
